\chapter{The Manufacturing Doctrine}
\label{ch:doctrine}

The manufacturing doctrine is a development discipline, not a testing strategy. It governs
how \texttt{wasm4pm-compat} is built: every commit is a receipt, every compile-fail fixture
is a proof obligation, and the commit history is itself an artifact whose certification state
can be measured exactly. This chapter formalizes the doctrine, defines the ALIVE gate that
certifies it, describes the PARTIAL checkpoint pattern that names residuals without false
claims, and documents the audit mesh that prevents regression.

%% ------------------------------------------------------------
\section{Commits as Receipts}
\label{sec:doctrine-receipts}

The central insight of the manufacturing doctrine is that a commit is not a unit of work. It
is a \emph{manufacturing transition}: a typed step from one certification state to the next,
backed by verifiable evidence of a specific law surface change.

\begin{definition}[Manufacturing Transition]
\label{def:manufacturing-transition}
A \emph{manufacturing transition} $\tau$ is a pair $(c, \rho)$ where $c$ is a git commit and
$\rho \in \mathcal{R}$ is a receipt class drawn from the set
\[
  \mathcal{R} = \{
    \texttt{paper-ledger},\;
    \texttt{paper-law},\;
    \texttt{type-law},\;
    \texttt{fixture-pass},\;
    \texttt{fixture-fail},\;
    \texttt{stderr},\;
    \texttt{ledger},\;
    \texttt{audit},\;
    \texttt{docs-law},\;
    \texttt{checkpoint},\;
    \texttt{tag}
  \}
\]
such that $c$ produces verifiable evidence of a specific law surface change traceable to
a paper, standard, or structural law in the process mining literature.
\end{definition}

Each receipt class has a distinct evidentiary role. Table~\ref{tab:receipt-classes} describes
the role, required artifact, and traceability requirement for each class.

\begin{table}[ht]
\centering
\caption{Manufacturing receipt classes: role, artifact, and traceability}
\label{tab:receipt-classes}
\begin{tabular}{@{}lll@{}}
\toprule
\textbf{Class} & \textbf{Artifact produced} & \textbf{Traceability requirement} \\
\midrule
\texttt{paper-ledger}  & Ledger row in \texttt{PAPER\_COVERAGE\_LEDGER.md} & Paper family + citation \\
\texttt{paper-law}     & Type surface in \texttt{src/*.rs} referencing the paper & Paper citation in rustdoc \\
\texttt{type-law}      & Named type or const-generic bound in \texttt{src/} & Named law from paper \\
\texttt{fixture-fail}  & Compile-fail \texttt{.rs} in \texttt{tests/ui/compile\_fail/} & Named violated law \\
\texttt{fixture-pass}  & Compile-pass \texttt{.rs} in \texttt{tests/ui/compile\_pass/} & Lawful path being proved \\
\texttt{stderr}        & \texttt{.stderr} snapshot paired with compile-fail fixture & Exact diagnostic text \\
\texttt{ledger}        & Updated ledger document & Commit hash + counts \\
\texttt{audit}         & Audit script in \texttt{scripts/audit/} & Gate criterion it checks \\
\texttt{docs-law}      & Law document in \texttt{docs/} & Law it specifies \\
\texttt{checkpoint}    & Checkpoint report in \texttt{docs/} & Gate state at time of commit \\
\texttt{tag}           & Git annotated tag & Certification level \\
\bottomrule
\end{tabular}
\end{table}

The receipt classes form a partial dependency order. A \texttt{fixture-fail} commit requires
a prior \texttt{type-law} or \texttt{paper-law} commit establishing the law being rejected.
A \texttt{stderr} commit is valid only if it is paired with the \texttt{fixture-fail} commit
that immediately precedes it. A \texttt{checkpoint} commit is valid only if the gate
measurements it records match the actual repository state at the time of the commit.

\begin{remark}[Receipt validity is structural, not procedural]
A commit that carries the correct receipt class prefix but does not produce the claimed
evidence is not a valid receipt. The manufacturing doctrine rejects procedural compliance:
a \texttt{fixture-fail} commit that introduces a fixture failing for the wrong reason
(a missing import, a typo, an unstable feature regression) is a defect, not a receipt. The
fixture must fail for the \emph{intended named law}, as verified by the exact
\texttt{.stderr} diagnostic.
\end{remark}

Reaching \textsc{PaperLaw Crown Alive~004} required 601 receipt-bearing commits manufactured
across May~2026. Table~\ref{tab:receipt-counts} shows the distribution across receipt classes.

\begin{table}[ht]
\centering
\caption{Receipt class distribution across 601 commits in the May~2026 manufacturing run}
\label{tab:receipt-counts}
\begin{tabular}{@{}lr@{}}
\toprule
\textbf{Receipt class} & \textbf{Commit count} \\
\midrule
\texttt{type-law}      & 140 \\
\texttt{fixture-pass}  &  89 \\
\texttt{paper-ledger}  &  68 \\
\texttt{fixture-fail}  &  64 \\
\texttt{paper-law}     &  63 \\
\texttt{audit}         &  38 \\
\texttt{docs-law}      &  18 \\
\texttt{checkpoint}    &  15 \\
\texttt{ledger}        &  11 \\
\texttt{stderr}        &  10 \\
Other (docs, chore)    &  85 \\
\midrule
\textbf{Total}         & \textbf{601} \\
\bottomrule
\end{tabular}
\end{table}

%% ------------------------------------------------------------
\section{The ALIVE Gate}
\label{sec:doctrine-alive}

The ALIVE gate is the certification surface for the manufacturing doctrine. A codebase
is \emph{ALIVE} at a given level if and only if a specific set of measurable gate criteria
are simultaneously satisfied on the same commit. The gate is not a narrative claim: every
criterion is a shell command that exits~0 on pass and exits non-zero on failure.

\begin{definition}[ALIVE Certification]
\label{def:alive}
Let $\mathcal{R}$ denote the repository state at a given commit. The ALIVE certification
$\text{ALIVE}_{004}(\mathcal{R})$ holds if and only if all ten gate criteria are satisfied:
\begin{equation}
\label{eq:alive-gate}
\text{ALIVE}_{004}(\mathcal{R}) \iff \bigwedge_{i=1}^{10} G_i(\mathcal{R})
\end{equation}
where each $G_i(\mathcal{R})$ is a measurable, machine-verifiable predicate on $\mathcal{R}$.
\end{definition}

The ten gate criteria for \textsc{PaperLaw Crown Alive~004} are:

\begin{enumerate}
  \item[\textbf{G1.}] \textbf{Paper Coverage.}
    $|\text{PAPER\_COVERAGE\_LEDGER}| \geq 80$ paper families, each with a registered
    type-law surface. Measurement: row count in \texttt{docs/PAPER\_COVERAGE\_LEDGER.md}.

  \item[\textbf{G2.}] \textbf{No Missing Type Law.}
    Zero modules have a paper registration without a corresponding type-law surface.
    $|\{m \in \text{modules} : \text{MISSING\_TYPE\_LAW}(m)\}| = 0$.
    Measurement: \texttt{bash scripts/audit/audit\_missing\_type\_law.sh}.

  \item[\textbf{G3.}] \textbf{Compile-Pass Fixtures.}
    $|\text{compile\_pass}| \geq 200$.
    Measurement: \texttt{ls tests/ui/compile\_pass/*.rs | wc -l}.

  \item[\textbf{G4.}] \textbf{Compile-Fail Fixtures.}
    $|\text{compile\_fail}| \geq 160$.
    Measurement: \texttt{ls tests/ui/compile\_fail/*.rs | wc -l}.

  \item[\textbf{G5.}] \textbf{Stderr Parity.}
    Every compile-fail fixture has a corresponding \texttt{.stderr} file with the expected
    compiler diagnostic. $|\text{compile\_fail}| = |\text{.stderr}|$.
    Measurement: \texttt{bash scripts/audit/audit\_stderr\_quality.sh}.

  \item[\textbf{G6.}] \textbf{Audit Script Coverage.}
    $|\text{audit scripts}| \geq 20$.
    Measurement: \texttt{ls scripts/audit/*.sh | wc -l}.

  \item[\textbf{G7.}] \textbf{Master Audit Gate.}
    All individual audit scripts pass.
    $\forall s \in \text{scripts/audit/}: \text{exit}(s) = 0$.
    Measurement: \texttt{bash scripts/audit/audit\_crown\_gate\_all.sh}.

  \item[\textbf{G8.}] \textbf{Cargo Tests.}
    $\text{exit}(\texttt{cargo test --all-features --tests}) = 0$.

  \item[\textbf{G9.}] \textbf{Clippy.}
    $\text{exit}(\texttt{cargo clippy --all-features -- -D warnings}) = 0$.
    No warnings are tolerated.

  \item[\textbf{G10.}] \textbf{Format.}
    $\text{exit}(\texttt{cargo fmt --check}) = 0$.
    No formatting deviations.
\end{enumerate}

\begin{remark}[Gate failure is a defect, not a discrepancy]
Following the Van der Aalst Constitution~\cite{vanderaalst2016process}: if any gate fails,
that is a defect, not a discrepancy. The crown cannot be awarded with exceptions, waivers,
or partial passes. A gate that passes most of the time is not a gate; it is a suggestion.
\end{remark}

The \textsc{PaperLaw Crown Alive~004} gate was passed on commit tagged
\texttt{wasm4pm-compat-paperlaw-crown-alive-004} with the following measured values:

\begin{table}[ht]
\centering
\caption{Measured gate values at \textsc{Crown Alive~004} certification}
\label{tab:crown-measurements}
\begin{tabular}{@{}lrr@{}}
\toprule
\textbf{Gate} & \textbf{Threshold} & \textbf{Measured} \\
\midrule
G1: Paper families       & $\geq 80$  & 98  \\
G2: Missing type law     & $= 0$      & 0   \\
G3: Compile-pass         & $\geq 200$ & 406 \\
G4: Compile-fail         & $\geq 160$ & 196 \\
G5: Stderr parity        & exact      & 196 \\
G6: Audit scripts        & $\geq 20$  & 23  \\
G7: Master audit gate    & exit 0     & PASS \\
G8: Cargo tests          & exit 0     & PASS \\
G9: Clippy               & exit 0     & PASS \\
G10: Format              & exit 0     & PASS \\
\bottomrule
\end{tabular}
\end{table}

%% ------------------------------------------------------------
\section{The PARTIAL Checkpoint Pattern}
\label{sec:doctrine-partial}

PARTIAL is not a failure state. It is a precise, honest state: the codebase satisfies some
but not all of the ALIVE gate criteria, and the residuals are named explicitly.

\begin{definition}[PARTIAL Checkpoint]
\label{def:partial}
A \emph{PARTIAL checkpoint} is a tagged commit that records:
\begin{enumerate}
  \item The set of gate criteria currently satisfied: $\{i : G_i(\mathcal{R}) = \top\}$.
  \item The set of residual criteria not yet satisfied: $\{i : G_i(\mathcal{R}) = \bot\}$.
  \item For each residual, a \emph{bill of materials}: the specific commits required to
        close it.
\end{enumerate}
\end{definition}

The bill-of-materials pattern is the operational core of PARTIAL. Rather than treating
unmet criteria as abstract shortfalls, the PARTIAL checkpoint enumerates them as a
manufactureable backlog. Each residual is expressed as one or more manufacturing
transitions of a specific receipt class:

\begin{example}[PARTIAL bill of materials]
\label{ex:partial-bom}
Suppose at checkpoint $\mathcal{R}_k$ that $G_3$ (compile-pass count) fails because
$|\text{compile\_pass}| = 185 < 200$. The bill of materials for residual $G_3$ is:
\[
  \text{BOM}(G_3) = \{15 \text{ commits of class } \texttt{fixture-pass}\}
\]
Each bill-of-materials entry names the paper family or module that the fixture must
address. A \texttt{fixture-pass} commit that does not reference a specific paper or
module is not a valid closure for $G_3$.
\end{example}

The PARTIAL checkpoint workflow is:

\[
  \text{PARTIAL} \;\to\; \text{residual inventory} \;\to\; \text{targeted closure}
  \;\to\; \text{recomputed gate} \;\to\; \text{ALIVE} \;\text{or named PARTIAL}
\]

The workflow terminates at ALIVE when all ten gates pass, or at a named PARTIAL when the
manufacturing work for a given sprint is exhausted. PARTIAL is better than false ALIVE
because:

\begin{enumerate}
  \item PARTIAL names exactly what is missing. ALIVE conceals it.
  \item PARTIAL enables targeted manufacturing. ALIVE produces undirected work.
  \item PARTIAL is auditable. False ALIVE requires defending an indefensible claim.
\end{enumerate}

Table~\ref{tab:alive-progression} shows how the PAPERLAW certification levels progressed
from \textsc{Alive~001} through \textsc{Crown Alive~004}, with the named PARTIAL states
between them.

\begin{table}[ht]
\centering
\caption{PAPERLAW milestone progression from \textsc{Alive~001} to \textsc{Crown Alive~004}}
\label{tab:alive-progression}
\begin{tabular}{@{}lrrrrl@{}}
\toprule
\textbf{Milestone} & \textbf{Fail} & \textbf{Pass} & \textbf{Papers} & \textbf{Modules} & \textbf{Status} \\
\midrule
\textsc{TypeLaw Alive~001}    &   4 &  10 &  8  & 10 & ALIVE \\
\textsc{PaperLaw Alive~002}   &  16 &  30 & 20  & 18 & ALIVE \\
\textsc{PaperLaw Partial~003} &  45 &  80 & 42  & 24 & PARTIAL (G1, G3, G4) \\
\textsc{PaperLaw Alive~003}   &  75 & 120 & 55  & 28 & ALIVE \\
\textsc{Crown Partial~004}    & 140 & 280 & 72  & 32 & PARTIAL (G3, G4) \\
\textsc{Crown Alive~004}      & 196 & 406 & 98  & 37 & ALIVE (Crown) \\
\bottomrule
\end{tabular}
\end{table}

%% ------------------------------------------------------------
\section{Anti-Regression Audit Mesh}
\label{sec:doctrine-audit}

The ALIVE gate measures forward progress. The audit mesh prevents backward regression.
It consists of 23~independent audit scripts, each checking a specific structural invariant
of the repository. Every script exits~0 on pass and exits~1 on fail, printing a named
defect when it fails.

The 23 scripts form two tiers:

\begin{description}
  \item[Hard audits] Structural invariants that, if violated, constitute immediate defects
    requiring remediation before any gate measurement. Hard audits check invariants that
    are unconditional: no missing type law, no engine creep, no stable-language claims,
    no unsealed refusals, no algorithm exports.

  \item[Soft audits] Coverage thresholds and quality checks where a single violation is
    a warning but the overall pass/fail depends on aggregate counts. Soft audits check
    fixture parity, benchmark coverage, and documentation completeness.
\end{description}

Table~\ref{tab:audit-scripts} lists all 23 audit scripts with their tier and invariant.

\begin{table}[ht]
\centering
\caption{Anti-regression audit mesh: 23 scripts, two tiers}
\label{tab:audit-scripts}
\begin{tabular}{@{}llp{6cm}@{}}
\toprule
\textbf{Script} & \textbf{Tier} & \textbf{Invariant enforced} \\
\midrule
\texttt{audit\_bench\_coverage}       & Soft & Benchmark suite covers all zero-cost types \\
\texttt{audit\_crosswalk\_links}      & Soft & Crosswalk links in \texttt{TYPE\_LAW\_CROSSWALK.md} resolve \\
\texttt{audit\_crown\_gate\_all}      & Hard & Orchestrator: runs all others, fails on any failure \\
\texttt{audit\_doctest\_disabled}     & Hard & \texttt{doctest = false} in \texttt{Cargo.toml} \\
\texttt{audit\_features}              & Hard & Exactly \{formats, strict, wasm4pm\}; no \texttt{nightly} feature \\
\texttt{audit\_graduation\_boundaries}& Hard & No engine types leak past \texttt{wasm4pm} feature gate \\
\texttt{audit\_id\_newtypes}          & Hard & All typed identifiers are newtypes, not bare aliases \\
\texttt{audit\_law\_annotations}      & Hard & Every public type has a law citation in its rustdoc \\
\texttt{audit\_metric\_bounds}        & Hard & All metric types carry \texttt{Between01} bounds \\
\texttt{audit\_missing\_type\_law}    & Hard & Zero modules have \texttt{MISSING\_TYPE\_LAW} in ledger \\
\texttt{audit\_module\_docs}          & Soft & Every module has a \texttt{//!} doc comment \\
\texttt{audit\_need9\_law}            & Hard & 8-bit state; Need9=split is rejected at the type level \\
\texttt{audit\_no\_algorithm\_exports}& Hard & No discovery/replay/alignment functions in \texttt{src/} \\
\texttt{audit\_no\_engine\_creep}     & Hard & No engine symbols (\texttt{discover}, \texttt{replay\_log}, \ldots) in \texttt{src/} \\
\texttt{audit\_no\_stable\_language}  & Hard & No affirmative stable-language claims in \texttt{src/} or \texttt{docs/} \\
\texttt{audit\_no\_unsealed\_refusal} & Hard & All refusal variants name a specific law (no \texttt{InvalidInput}) \\
\texttt{audit\_paper\_law\_ledger}    & Hard & Every ledger row has a valid classification status \\
\texttt{audit\_pass\_fail\_pairs}     & Soft & Each compile-fail law has a corresponding compile-pass law \\
\texttt{audit\_projection\_loss}      & Hard & All projections carry \texttt{LossPolicy} + \texttt{LossReport} \\
\texttt{audit\_receipt\_chain}        & Hard & Receipt chain types are structurally sound \\
\texttt{audit\_stderr\_quality}       & Hard & Every compile-fail \texttt{.rs} has an exact-match \texttt{.stderr} \\
\texttt{audit\_trybuild\_receipts}    & Hard & All compile-fail \texttt{.rs} files compile-fail for the named law \\
\texttt{audit\_witness\_markers}      & Hard & All witness types implement the \texttt{Witness} sealed trait \\
\bottomrule
\end{tabular}
\end{table}

The orchestrator script \texttt{audit\_crown\_gate\_all.sh} runs all 22 subordinate scripts
in sequence, reports pass/fail per script, and exits~0 only if all pass. This makes Gate~7
of the ALIVE certification ($G_7$ in \S\ref{sec:doctrine-alive}) a single atomic check.

Algorithm~\ref{alg:crown-audit} formalizes the crown audit protocol.

\begin{algorithm}[ht]
\caption{Crown Audit Protocol (\texttt{audit\_crown\_gate\_all.sh})}
\label{alg:crown-audit}
\DontPrintSemicolon
\KwIn{Repository state $\mathcal{R}$, audit script set $\mathcal{S} = \{s_1, \ldots, s_{22}\}$}
\KwOut{PASS (exit 0) or FAIL (exit 1) with named defects}

$\text{PASS} \leftarrow 0$\;
$\text{FAIL} \leftarrow 0$\;
$\text{defects} \leftarrow \emptyset$\;

\ForEach{$s_i \in \mathcal{S} \setminus \{\texttt{audit\_crown\_gate\_all}\}$}{
  \eIf{$\text{exit}(s_i(\mathcal{R})) = 0$}{
    $\text{PASS} \leftarrow \text{PASS} + 1$\;
    \textbf{print} \texttt{PASS} $s_i$\;
  }{
    $\text{FAIL} \leftarrow \text{FAIL} + 1$\;
    $\text{defects} \leftarrow \text{defects} \cup \{s_i\}$\;
    \textbf{print} \texttt{FAIL} $s_i$\;
  }
}

\textbf{print} \texttt{Crown Audit Gate:} $\text{PASS}$ \texttt{pass,} $\text{FAIL}$ \texttt{fail}\;

\eIf{$\text{FAIL} = 0$}{
  \Return{exit~0}\;
}{
  \ForEach{$d \in \text{defects}$}{
    \textbf{print} \texttt{DEFECT:} $d$\;
  }
  \Return{exit~1}\;
}
\end{algorithm}

A key design property of the audit mesh is monotonicity: adding new audit scripts can only
raise the bar for future ALIVE certifications. Removing a script requires an explicit
defect record explaining why the invariant it checked is no longer relevant. No script has
been removed since \textsc{Alive~002}.

%% ------------------------------------------------------------
\section{Velocity and Scale}
\label{sec:doctrine-velocity}

The manufacturing doctrine is not a theoretical framework. It was exercised at scale during
May~2026, producing 601 receipt-bearing commits that advanced the codebase from
\textsc{PaperLaw Alive~002} (16 compile-fail fixtures, 30 compile-pass fixtures, 20 papers)
to \textsc{Crown Alive~004} (196 compile-fail receipts, 406 compile-pass receipts, 98
papers).

This velocity was enabled by workflow orchestration using Claude Code dynamic workflows.
The orchestration pattern that made it possible was the \emph{schema-free agent pattern}.

Earlier orchestration attempts used structured output schemas to route agent subtasks.
Under high parallelism, the structured output path triggered a \texttt{StructuredOutput}
timeout in the tool infrastructure: the schema validation for large nested response objects
exceeded the maximum allowed latency, causing entire agent batches to time out and leaving
partial work uncommitted.

The schema-free agent pattern resolves this by returning all agent results as plain text
with a documented parsing convention. The orchestrating workflow applies a lightweight
text-pattern extractor rather than a schema validator. The manufacturing pipeline then
routes the extracted result to the appropriate receipt class. This change reduced per-agent
latency from timeout-frequent to sub-second in the common case, enabling the parallel
manufacturing of fixture pairs (compile-fail + stderr) in a single orchestration cycle.

The broader lesson is that manufacturing velocity in a type-law codebase is bottlenecked
not by the number of available law surfaces (which is bounded by the paper corpus) but by
the throughput of the orchestration layer. The schema-free agent pattern unlocked a
throughput regime that made 601 commits manufacturable within a single calendar month
without sacrificing receipt validity.

The manufacturing pipeline at full throughput operates as follows:

\begin{enumerate}
  \item \textbf{Paper intake.} The paper corpus is scanned for unregistered families. Each
        unregistered family generates a \texttt{paper-ledger} task.

  \item \textbf{Type-law surface assignment.} Each ledgered paper is assigned to a module.
        The module receives a \texttt{paper-law} task that produces a Rust type surface.

  \item \textbf{Fixture pair manufacturing.} Each new type-law surface generates two tasks:
        a \texttt{fixture-fail} task (compile-fail fixture + stderr receipt) and a
        \texttt{fixture-pass} task (compile-pass fixture proving the lawful path is open).

  \item \textbf{Audit gate.} After each batch of fixture pairs, the crown audit runner
        checks all 23 audit scripts. Any failure generates a targeted \texttt{audit} or
        \texttt{type-law} remediation task.

  \item \textbf{Checkpoint.} When a gate stratum is reached (e.g., all G3--G6 criteria
        pass), a \texttt{checkpoint} commit records the state. If all ten gates pass, a
        \texttt{tag} commit awards the ALIVE level.
\end{enumerate}

%% ------------------------------------------------------------
\section{Paper as Law}
\label{sec:doctrine-paper}

The PAPERLAW series connects the manufacturing doctrine to the process mining literature.
Every paper in the corpus is classified under one of five statuses that determine what type
coverage it requires.

\begin{description}
  \item[\texttt{COVERED\_BY\_TYPE}] The paper introduces formal objects (types, structures,
    laws) that are directly encoded as zero-cost Rust types in this crate. The structural
    shapes are present; execution logic is absent by design.

  \item[\texttt{COVERED\_BY\_GRADUATION\_BOUNDARY}] The paper's primary contribution is an
    algorithm (discovery, conformance checking, replay, alignment, prediction, query
    evaluation). The algorithm's output shapes may be typed here; the algorithm itself
    graduates to \texttt{wasm4pm}. The ledger entry names what specifically graduates using
    a \texttt{GraduationReason} variant.

  \item[\texttt{PARTIAL\_WITH\_REASON}] Some of the paper's formal objects are typed here;
    others are missing or incorrectly classified. The ledger entry names what is absent and
    why. This is a legitimate state, not a failure.

  \item[\texttt{DUPLICATE\_OR\_BACKGROUND}] The paper repeats or summarises content from
    another ledgered paper, or provides background context without introducing new formal
    objects. The ledger entry names the primary citation.

  \item[\texttt{OUT\_OF\_SCOPE\_WITH\_REASON}] The paper does not introduce process-mining
    formal objects. The reason must be specific. A paper that introduces any formal object
    that could be a zero-cost Rust type cannot be \texttt{OUT\_OF\_SCOPE}.
\end{description}

\begin{remark}[\texttt{MISSING\_TYPE\_LAW} is the forbidden status]
\label{rem:missing-type-law}
\texttt{MISSING\_TYPE\_LAW} is not a valid classification status. A paper either has
coverage, a named graduation boundary, a documented partial, or a specific out-of-scope
reason. Hidden gaps are defects. Gate~G2 of the ALIVE certification enforces this: the
audit script \texttt{audit\_missing\_type\_law.sh} exits non-zero if any ledger row
carries \texttt{MISSING\_TYPE\_LAW}. No path to ALIVE is open while this defect exists.
\end{remark}

The relationship between paper classification and type coverage is one-to-one: every
\texttt{COVERED\_BY\_TYPE} row in the ledger must correspond to a named type surface in
\texttt{src/*.rs} with a rustdoc citation referencing the paper. The audit script
\texttt{audit\_law\_annotations.sh} verifies this correspondence and fails if any
\texttt{COVERED\_BY\_TYPE} paper lacks a law annotation in source.

Table~\ref{tab:paperlaw-milestones} shows the full PAPERLAW milestone progression with
fixture counts, paper coverage, and module counts at each certification level.

\begin{table}[ht]
\centering
\caption{PAPERLAW milestone progression: fixture counts, papers, and modules}
\label{tab:paperlaw-milestones}
\begin{tabular}{@{}lrrrrrl@{}}
\toprule
\textbf{Milestone} & \textbf{Fail} & \textbf{Pass} & \textbf{.stderr} & \textbf{Papers} & \textbf{Modules} & \textbf{Verdict} \\
\midrule
\textsc{TypeLaw Alive~001}
  &   4 &  10 &   4 &  8  & 10 & ALIVE \\
\textsc{PaperLaw Alive~002}
  &  16 &  30 &  16 & 20  & 18 & ALIVE \\
\textsc{PaperLaw Partial~003}
  &  45 &  80 &  45 & 42  & 24 & PARTIAL \\
\textsc{PaperLaw Alive~003}
  &  75 & 120 &  75 & 55  & 28 & ALIVE \\
\textsc{Crown Partial~004}
  & 140 & 280 & 140 & 72  & 32 & PARTIAL \\
\textsc{Crown Alive~004}
  & 196 & 406 & 196 & 98  & 37 & ALIVE (Crown) \\
\bottomrule
\end{tabular}
\end{table}

The progression from \textsc{TypeLaw Alive~001} to \textsc{Crown Alive~004} represents a
$49\times$ increase in compile-fail fixtures, a $40\times$ increase in compile-pass
fixtures, and a $12\times$ increase in paper coverage. The \texttt{.stderr} receipt column
maintains exact parity with compile-fail throughout: at no certification level is there a
compile-fail fixture without a corresponding \texttt{.stderr} diagnostic snapshot.

This parity is not incidental. It is enforced by Gate~G5 and by the \texttt{audit\_stderr\_quality.sh}
hard audit. A compile-fail fixture without an exact \texttt{.stderr} receipt is not a valid
type-law receipt: it cannot prove that the fixture is failing for the intended named law
rather than for an accidental reason.

The manufacturing doctrine closes with a statement of its invariant:

\begin{quote}
Start with paper law. Admit through types. Prove through receipts. Graduate to execution.
\end{quote}

Every manufacturing transition either advances this trajectory or it does not count. The
commit history is the proof. The ALIVE gate is the judge.
